% !TeX root = ../main.tex
\documentclass[../main.tex]{subfiles}

\begin{document}
\chapter{Potenciales retardados y fuentes en movimiento}
\label{sec:potenciales-retardados-fuentes-movimiento}

En el gauge de Lorenz, las ecuaciones de Maxwell inhomogéneas se transforman en ecuaciones de onda inhomogéneas para los potenciales. Esta observación es el punto de partida natural para discutir \textbf{potenciales retardados}: si las fuentes cambian en el tiempo, el potencial observado en un punto no puede depender de la fuente ``en el mismo instante'' de manera instantánea, sino de la fuente evaluada en el tiempo compatible con la propagación a rapidez $c$.

En esta sección seguimos la estructura del capítulo de radiación de las notas del curso, pero la incorporamos en la notación y el estilo del presente apunte. La idea física central es sencilla: resolver las ecuaciones de onda con fuente equivale a reconstruir el campo en un evento $(ct,\vec x)$ a partir de la información de las fuentes en su \textbf{cono de luz pasado}.

\section{Ecuaciones de onda inhomogéneas para los potenciales}
\label{subsec:onda-inhomogenea-potenciales}

Recordemos que, al imponer el gauge de Lorenz,
\begin{equation*}
    \vec\nabla\cdot\vec A+\frac{1}{c^2}\frac{\partial\phi}{\partial t}=0,
\end{equation*}
las ecuaciones para los potenciales se desacoplan como
\begin{align}
    \nabla^2\phi-\frac{1}{c^2}\frac{\partial^2\phi}{\partial t^2}
    &=-\frac{\rho}{\varepsilon},
    \label{eq:rad-onda-phi}\\
    \nabla^2\vec A-\frac{1}{c^2}\frac{\partial^2\vec A}{\partial t^2}
    &=-\mu\vec J.
    \label{eq:rad-onda-A}
\end{align}
Aquí $\varepsilon$ y $\mu$ son las constantes del medio homogéneo considerado, y $c=1/\sqrt{\varepsilon\mu}$. En vacío se reemplazan por $\varepsilon_0$ y $\mu_0$.

Conviene escribir ambas ecuaciones de manera unificada. Definimos el operador de onda
\begin{equation}
    \Box:=\nabla^2-\frac{1}{c^2}\frac{\partial^2}{\partial t^2}.
    \label{eq:rad-box-def}
\end{equation}
Entonces el problema que debemos resolver tiene la forma general
\begin{equation}
    \Box\Psi(\vec x,t)=-g(\vec x,t),
    \label{eq:rad-onda-inhomogenea-general}
\end{equation}
donde $\Psi$ representa una componente del potencial escalar o vectorial, y $g$ representa la fuente correspondiente.

\begin{cajaidea}
\textbf{Idea física.} Las ecuaciones \eqref{eq:rad-onda-phi} y \eqref{eq:rad-onda-A} no son ecuaciones de Poisson instantáneas, sino ecuaciones de onda con fuente. Por tanto, la solución debe recordar que la información electromagnética se propaga con rapidez finita $c$.
\end{cajaidea}

\section{Reducción temporal a una ecuación de Helmholtz}
\label{subsec:helmholtz-green-retardados}

Para resolver \eqref{eq:rad-onda-inhomogenea-general}, introducimos la transformada de Fourier temporal
\begin{equation}
    \widetilde{\Psi}(\vec x,\omega):=
    \int_{-\infty}^{\infty}\Psi(\vec x,t)\,\ee^{-\ii\omega t}\,dt,
    \label{eq:rad-fourier-psi}
\end{equation}
\begin{equation}
    \widetilde g(\vec x,\omega):=
    \int_{-\infty}^{\infty}g(\vec x,t)\,\ee^{-\ii\omega t}\,dt.
    \label{eq:rad-fourier-g}
\end{equation}
La transformada inversa queda dada por
\begin{equation}
    \Psi(\vec x,t)=\frac{1}{2\pi}\int_{-\infty}^{\infty}
    \widetilde\Psi(\vec x,\omega)\,\ee^{\ii\omega t}\,d\omega,
    \qquad
    g(\vec x,t)=\frac{1}{2\pi}\int_{-\infty}^{\infty}
    \widetilde g(\vec x,\omega)\,\ee^{\ii\omega t}\,d\omega.
    \label{eq:rad-fourier-inversa}
\end{equation}
Al transformar la ecuación de onda, la derivada temporal de segundo orden produce un factor $-\omega^2$. Por tanto,
\begin{equation*}
    \mathcal F_t\left[\frac{\partial^2\Psi}{\partial t^2}\right]
    =-\omega^2\widetilde\Psi,
\end{equation*}
y \eqref{eq:rad-onda-inhomogenea-general} se convierte en
\begin{equation}
    \left(\nabla^2+k^2\right)\widetilde\Psi(\vec x,\omega)
    =-\widetilde g(\vec x,\omega),
    \qquad
    k:=\frac{\omega}{c}.
    \label{eq:rad-helmholtz-inhomogenea}
\end{equation}
Esta es una ecuación de Helmholtz inhomogénea para cada frecuencia $\omega$.

\section{Función de Green independiente del tiempo}
\label{subsec:green-helmholtz-retardada}

Para resolver \eqref{eq:rad-helmholtz-inhomogenea}, buscamos una función de Green $G(\vec x,\vec x')$ que satisfaga
\begin{equation}
    \left(\nabla^2+k^2\right)G(\vec x,\vec x')
    =-\delta^{(3)}(\vec x-\vec x').
    \label{eq:rad-green-helmholtz}
\end{equation}
Con esta convención, una solución particular queda escrita como
\begin{equation}
    \widetilde\Psi(\vec x,\omega)=
    \int G(\vec x,\vec x')\,\widetilde g(\vec x',\omega)\,dV'.
    \label{eq:rad-solucion-green-frecuencia}
\end{equation}
La solución esférica saliente o entrante de \eqref{eq:rad-green-helmholtz} es
\begin{equation}
    G_{\pm}(\vec x,\vec x')=
    \frac{\ee^{\pm \ii k |\vec x-\vec x'|}}{4\pi |\vec x-\vec x'|}.
    \label{eq:rad-green-pm}
\end{equation}
Por tanto,
\begin{equation}
    \widetilde\Psi_{\pm}(\vec x,\omega)=
    \frac{1}{4\pi}\int
    \frac{\widetilde g(\vec x',\omega)\,\ee^{\pm \ii k |\vec x-\vec x'|}}
    {|\vec x-\vec x'|}\,dV'.
    \label{eq:rad-psi-frecuencia-pm}
\end{equation}
Al volver al dominio temporal, obtenemos
\begin{equation}
    \Psi_{\pm}(\vec x,t)=
    \frac{1}{4\pi}\int
    \frac{g(\vec x',t_{\pm})}{|\vec x-\vec x'|}\,dV',
    \label{eq:rad-solucion-pm-temporal}
\end{equation}
donde
\begin{equation}
    t_{\pm}:=t\pm\frac{|\vec x-\vec x'|}{c}.
    \label{eq:rad-tiempo-pm}
\end{equation}
La solución con $t_-=t-|\vec x-\vec x'|/c$ es la solución \textbf{retardada}; la solución con $t_+=t+|\vec x-\vec x'|/c$ es la solución \textbf{avanzada}.

\begin{cuidado}
La solución avanzada no es matemáticamente inválida; lo que ocurre es que no representa la elección causal usual. Para predecir los campos a partir de fuentes dadas, se escoge la solución retardada, porque sólo usa información situada en el pasado causal del evento de observación.
\end{cuidado}

\section{Potenciales retardados}
\label{subsec:potenciales-retardados}

Aplicando \eqref{eq:rad-solucion-pm-temporal} a las ecuaciones \eqref{eq:rad-onda-phi} y \eqref{eq:rad-onda-A}, los potenciales retardados quedan
\begin{equation}
    \boxed{
    \phi_{\mathrm{ret}}(\vec x,t)=
    \frac{1}{4\pi\varepsilon}\int
    \frac{\rho(\vec x',t_{\mathrm{ret}})}{|\vec x-\vec x'|}\,dV'
    }
    \label{eq:rad-potencial-escalar-retardado}
\end{equation}
\begin{equation}
    \boxed{
    \vec A_{\mathrm{ret}}(\vec x,t)=
    \frac{\mu}{4\pi}\int
    \frac{\vec J(\vec x',t_{\mathrm{ret}})}{|\vec x-\vec x'|}\,dV'
    }
    \label{eq:rad-potencial-vectorial-retardado}
\end{equation}
con
\begin{equation}
    \boxed{
    t_{\mathrm{ret}}(\vec x,\vec x',t)=
    t-\frac{|\vec x-\vec x'|}{c}.
    }
    \label{eq:rad-tiempo-retardado-distribucion}
\end{equation}
Estas expresiones son formalmente parecidas a los potenciales electrostático y magnetostático, pero con una diferencia esencial: las fuentes no se evalúan en $t$, sino en el tiempo retardado $t_{\mathrm{ret}}$.

La condición \eqref{eq:rad-tiempo-retardado-distribucion} puede escribirse como
\begin{equation}
    c^2(t-t_{\mathrm{ret}})^2-|\vec x-\vec x'|^2=0.
    \label{eq:rad-relacion-tipo-luz-distribucion}
\end{equation}
Por tanto, el evento fuente $(ct_{\mathrm{ret}},\vec x')$ y el evento de observación $(ct,\vec x)$ están conectados por una señal luminosa.

\begin{figure}[ht!]
\centering
\begin{tikzpicture}[scale=0.95]
    \draw[eje] (-0.2,0) -- (7.1,0) node[right] {$x$};
    \draw[eje] (0,-0.2) -- (0,4.7) node[above] {$ct$};

    \coordinate (O) at (4.8,3.7);
    \coordinate (S) at (2.2,1.1);

    \draw[UdeCAzul, thick] (O) -- (2.5,1.4);
    \draw[UdeCAzul, thick] (O) -- (7.0,1.5);
    \draw[UdeCAzul!45, dashed] (O) -- (3.3,0.2);
    \draw[UdeCAzul!45, dashed] (O) -- (6.3,0.2);

    \fill[UdeCDorado] (O) circle (2.2pt) node[above right] {$(ct,\vec x)$};
    \fill[UdeCAzul] (S) circle (2.2pt) node[below left] {$(ct_{\mathrm{ret}},\vec x')$};
    \draw[-{Latex}, thick, UdeCDorado] (S) -- node[above, sloped] {señal a rapidez $c$} (O);

    \node[align=center, text=UdeCAzul] at (5.85,2.1) {cono de luz\ pasado};
\end{tikzpicture}
\caption{Interpretación causal de los potenciales retardados: el evento fuente debe estar sobre el cono de luz pasado del evento de observación.}
\label{fig:rad-cono-luz-retardado}
\end{figure}

\begin{cajaidea}
\textbf{Lectura física.} Los potenciales retardados permiten predecir el campo en $(\vec x,t)$ conociendo las fuentes en tiempos anteriores. Un cambio en una distribución de carga o corriente no se transmite instantáneamente: su efecto llega después de un tiempo de vuelo $|\vec x-\vec x'|/c$.
\end{cajaidea}

\section{Ecuaciones de Jefimenko}
\label{subsec:ecuaciones-jefimenko}

A partir de los potenciales retardados se obtienen los campos mediante
\begin{equation*}
    \E=-\vec\nabla\phi-\frac{\partial\vec A}{\partial t},
    \qquad
    \B=\vec\nabla\times\vec A.
\end{equation*}
Al realizar cuidadosamente las derivadas, recordando que $t_{\mathrm{ret}}$ depende de $\vec x$, se llega a las \textbf{ecuaciones de Jefimenko}:
\begin{equation}
    \boxed{
    \E(\vec x,t)=\frac{1}{4\pi\varepsilon}\int
    \left[
    \frac{\rho(\vec x',t_{\mathrm{ret}})(\vec x-\vec x')}{|\vec x-\vec x'|^3}
    +\frac{1}{c}\frac{\dot\rho(\vec x',t_{\mathrm{ret}})(\vec x-\vec x')}{|\vec x-\vec x'|^2}
    -\frac{1}{c^2}\frac{\dot{\vec J}(\vec x',t_{\mathrm{ret}})}{|\vec x-\vec x'|}
    \right]dV'
    }
    \label{eq:rad-jefimenko-E}
\end{equation}
\begin{equation}
    \boxed{
    \B(\vec x,t)=\frac{\mu}{4\pi}\int
    \left[
    \frac{\vec J(\vec x',t_{\mathrm{ret}})\times(\vec x-\vec x')}{|\vec x-\vec x'|^3}
    +\frac{1}{c}\frac{\dot{\vec J}(\vec x',t_{\mathrm{ret}})\times(\vec x-\vec x')}{|\vec x-\vec x'|^2}
    \right]dV'
    }
    \label{eq:rad-jefimenko-B}
\end{equation}
Aquí el punto denota derivada parcial respecto de la variable temporal de la fuente; por ejemplo,
\begin{equation*}
    \dot\rho(\vec x',t_{\mathrm{ret}})=
    \left.\frac{\partial\rho(\vec x',\tau)}{\partial\tau}\right|_{\tau=t_{\mathrm{ret}}}.
\end{equation*}

Las ecuaciones de Jefimenko son una versión explícitamente causal de los campos. Si las fuentes son estacionarias, entonces $\dot\rho=0$ y $\dot{\vec J}=\vec 0$, recuperándose las expresiones estáticas usuales. En cambio, cuando las fuentes varían en el tiempo, aparecen términos adicionales que dependen de la variación temporal de la carga y de la corriente.

\section{Potenciales de Liénard--Wiechert}
\label{subsec:potenciales-lienard-wiechert}

Los \textbf{potenciales de Liénard--Wiechert} son los potenciales retardados producidos por una carga puntual $q$ que se mueve siguiendo una trayectoria conocida $\vec z(t)$. Es decir, ahora la fuente no está distribuida en el volumen, sino concentrada sobre la línea de mundo de la partícula.

Las densidades de carga y corriente asociadas a una carga puntual en movimiento son
\begin{equation}
    \rho(\vec x,t)=q\delta^{(3)}(\vec x-\vec z(t)),
    \qquad
    \vec J(\vec x,t)=q\vec v(t)\delta^{(3)}(\vec x-\vec z(t)),
    \label{eq:rad-densidades-carga-puntual}
\end{equation}
donde
\begin{equation*}
    \vec v(t)=\frac{d\vec z}{dt}.
\end{equation*}
Al reemplazar \eqref{eq:rad-densidades-carga-puntual} en \eqref{eq:rad-potencial-escalar-retardado} y \eqref{eq:rad-potencial-vectorial-retardado}, se obtiene
\begin{equation}
    \phi_{\mathrm{ret}}(\vec x,t)=
    \frac{q}{4\pi\varepsilon}\int
    \frac{\delta^{(3)}\!\bigl(\vec x'-\vec z(t-|\vec x-\vec x'|/c)\bigr)}{|\vec x-\vec x'|}\,d^3x',
    \label{eq:rad-lw-phi-integral-inicial}
\end{equation}
\begin{equation}
    \vec A_{\mathrm{ret}}(\vec x,t)=
    \frac{q\mu}{4\pi}\int
    \frac{\vec v(t-|\vec x-\vec x'|/c)\,
    \delta^{(3)}\!\bigl(\vec x'-\vec z(t-|\vec x-\vec x'|/c)\bigr)}{|\vec x-\vec x'|}\,d^3x'.
    \label{eq:rad-lw-A-integral-inicial}
\end{equation}
La dificultad está en que la delta de Dirac contiene a $\vec x'$ tanto de manera explícita como dentro del tiempo retardado. Para aislar correctamente la contribución de la trayectoria, consideramos una integral genérica de la forma
\begin{equation}
    I(\vec x,t)=\int
    f(\vec x,\vec x',t)\,
    \delta^{(3)}\!\bigl(\vec x'-\vec z(t-|\vec x-\vec x'|/c)\bigr)\,d^3x'.
    \label{eq:rad-lw-integral-generica}
\end{equation}
Introducimos una variable temporal auxiliar $T$ y escribimos
\begin{equation}
    \delta^{(3)}\!\bigl(\vec x'-\vec z(t-|\vec x-\vec x'|/c)\bigr)
    =\int_{-\infty}^{\infty}
    \delta^{(3)}(\vec x'-\vec z(T))
    \delta\!\bigl(T-t+|\vec x-\vec x'|/c\bigr)\,dT.
    \label{eq:rad-lw-delta-auxiliar}
\end{equation}
Entonces
\begin{align}
    I(\vec x,t)
    &=\int_{-\infty}^{\infty}\int
    f(\vec x,\vec x',t)\,
    \delta^{(3)}(\vec x'-\vec z(T))
    \delta\!\bigl(T-t+|\vec x-\vec x'|/c\bigr)
    \,d^3x'\,dT \\
    &=\int_{-\infty}^{\infty}
    f(\vec x,\vec z(T),t)\,
    \delta\!\bigl(T-t+|\vec x-\vec z(T)|/c\bigr)\,dT.
    \label{eq:rad-lw-integral-sobre-T}
\end{align}
Definimos
\begin{equation}
    g(T):=T-t+\frac{1}{c}|\vec x-\vec z(T)|.
    \label{eq:rad-lw-gT}
\end{equation}
El tiempo retardado $t'$ queda determinado por la raíz de $g$, es decir,
\begin{equation}
    \boxed{
    t'=t-\frac{|\vec x-\vec z(t')|}{c}.
    }
    \label{eq:rad-lw-tiempo-retardado}
\end{equation}
Esta es una ecuación implícita: para conocer el tiempo retardado hay que encontrar el punto de la trayectoria de la carga que se conecta causalmente con el evento de observación.

Usando la identidad de la delta de Dirac
\begin{equation*}
    \delta(g(T))=\sum_i\frac{\delta(T-T_i)}{|g'(T_i)|},
\end{equation*}
y suponiendo movimiento subluminal, de modo que existe una única raíz retardada relevante, se tiene
\begin{equation}
    I(\vec x,t)=
    \frac{f(\vec x,\vec z(t'),t)}{|g'(t')|}.
    \label{eq:rad-lw-I-con-gprima}
\end{equation}
Calculemos ahora $g'(T)$. Por regla de la cadena,
\begin{align}
    g'(T)
    &=1+\frac{1}{c}\frac{d}{dT}|\vec x-\vec z(T)| \\
    &=1+\frac{1}{c}\frac{\vec x-\vec z(T)}{|\vec x-\vec z(T)|}\cdot\frac{d}{dT}(\vec x-\vec z(T)) \\
    &=1-\frac{1}{c}\frac{\vec v(T)\cdot(\vec x-\vec z(T))}{|\vec x-\vec z(T)|}.
    \label{eq:rad-lw-gprima-calculo}
\end{align}
Definimos, evaluado en el tiempo retardado,
\begin{equation}
    \vec{\mathcal R}:=\vec x-\vec z(t'),
    \qquad
    \mathcal R:=|\vec x-\vec z(t')|,
    \qquad
    \hat n:=\frac{\vec{\mathcal R}}{\mathcal R},
    \qquad
    \vec\beta:=\frac{\vec v(t')}{c}.
    \label{eq:rad-lw-definiciones-retardadas}
\end{equation}
Con esta notación,
\begin{equation}
    g'(t')=1-\vec\beta\cdot\hat n.
    \label{eq:rad-lw-gprima-final}
\end{equation}
Como $|\vec\beta|<1$, se tiene $g'(t')>0$ para la raíz retardada, y por tanto
\begin{equation}
    I(\vec x,t)=
    \frac{f(\vec x,\vec z(t'),t)}{1-\vec\beta\cdot\hat n}.
    \label{eq:rad-lw-I-final}
\end{equation}

Aplicando este resultado a \eqref{eq:rad-lw-phi-integral-inicial} y \eqref{eq:rad-lw-A-integral-inicial}, obtenemos finalmente
\begin{equation}
    \boxed{
    \phi_{\mathrm{LW}}(\vec x,t)=
    \frac{q}{4\pi\varepsilon}
    \left.\frac{1}{\mathcal R(1-\vec\beta\cdot\hat n)}\right|_{\mathrm{ret}}
    }
    \label{eq:rad-lw-phi-final}
\end{equation}
\begin{equation}
    \boxed{
    \vec A_{\mathrm{LW}}(\vec x,t)=
    \frac{\mu c q}{4\pi}
    \left.\frac{\vec\beta}{\mathcal R(1-\vec\beta\cdot\hat n)}\right|_{\mathrm{ret}}
    }
    \label{eq:rad-lw-A-final}
\end{equation}
Estas son las expresiones de los potenciales de Liénard--Wiechert.

\begin{cajaidea}
\textbf{Resumen conceptual.} Para una carga puntual en movimiento arbitrario, los potenciales en $(\vec x,t)$ dependen de la posición y velocidad de la carga en el único instante retardado $t'$ que satisface \eqref{eq:rad-lw-tiempo-retardado}. La expresión no sólo contiene la distancia retardada $\mathcal R$, sino también el factor cinemático $1-\vec\beta\cdot\hat n$, que proviene del cambio de variable dentro de la delta de Dirac.
\end{cajaidea}

Además, usando $c^{-2}=\varepsilon\mu$, las expresiones \eqref{eq:rad-lw-phi-final} y \eqref{eq:rad-lw-A-final} implican la relación
\begin{equation}
    \boxed{
    \vec A_{\mathrm{LW}}(\vec x,t)=
    \frac{\phi_{\mathrm{LW}}(\vec x,t)}{c}\left.\vec\beta\right|_{\mathrm{ret}}
    =\frac{\phi_{\mathrm{LW}}(\vec x,t)}{c^2}\left.\vec v\right|_{\mathrm{ret}}.
    }
    \label{eq:rad-lw-relacion-A-phi}
\end{equation}

\begin{figure}[H]
\centering
\begin{tikzpicture}[scale=0.95]
    \draw[eje] (-0.2,0) -- (7.3,0) node[right] {$x$};
    \draw[eje] (0,-0.2) -- (0,4.8) node[above] {$ct$};

    \draw[UdeCAzul, very thick, domain=0.6:4.0, smooth, variable=\s]
        plot ({1.4+0.9*\s+0.25*sin(90*\s)}, {0.45+\s});

    \coordinate (P) at (5.6,4.0);
    \coordinate (Q) at (3.1,2.05);

    \fill[UdeCDorado] (P) circle (2.2pt) node[above right] {$(ct,\vec x)$};
    \fill[UdeCAzul] (Q) circle (2.2pt) node[below left] {$(ct',\vec z(t'))$};

    \draw[UdeCAzul!55, dashed, thick] (P) -- (3.2,1.6);
    \draw[UdeCAzul!55, dashed, thick] (P) -- (7.2,1.6);
    \draw[-{Latex}, thick, UdeCDorado] (Q) -- node[above, sloped] {$\mathcal R=c(t-t')$} (P);

    \node[align=center, text=UdeCAzul] at (2.35,4.35) {línea de mundo\ de la carga};
    \node[align=center, text=UdeCAzul] at (6.35,2.25) {cono de luz\ pasado};
\end{tikzpicture}
\caption{El tiempo retardado $t'$ se obtiene como la intersección entre la línea de mundo de la carga y el cono de luz pasado del evento de observación.}
\label{fig:rad-lienard-wiechert-retardo}
\end{figure}

\end{document}
